\chapter{كيف نصلح آلة بلا مخطط}
% full version: chapters/17_debugging.tex

\pencilfig{17}{\pencilart{17}}{لوحة بلا مخطط: المسبار لا يسمع إلا بضعة أسلاك من مليارات.}

تخيّل أنهم أعطوك لوحة إلكترونية تعمل بلا مخطط وطلبوا منك إصلاحها. لدى المهندس أربع أدوات: مسبار ليصغي إلى الإشارة؛ ومولّد ليحقن إشارته الخاصة؛ وإمكانية فصل قطعة لرؤية ما يتعطل؛ والوصول إلى سجلات التحكم. وبهذه الأدوات بالضبط يُدرس الدماغ أيضًا.

\section*{مسبار لألف عصبون}

القطب المغروس في الدماغ يسمع نقرات بضعة عصبونات قريبة. قد يبدو أن ألف قطب لا شيء من ستة وثمانين مليارًا. لكنها تكفي للتحكم في مؤشر على الشاشة. لماذا؟ العصبونات المتجاورة تضج بطريقة متشابهة، والنشاط الذي يتحكم في اليد يوصف في الحقيقة بعشرة أو عشرين متغيرًا مشتركًا فقط. لا داعي للإصغاء إلى الجميع، يكفي سماع هذه ”الأصوات“ القليلة.

الإشارة الخام من ألف قطب مئات الملايين من البتات في الثانية: قناة لاسلكية من تحت الجمجمة لن تمررها. أما النقرات نفسها فتحمل نحو ألف مرة أقل. لذلك يحسن استخلاص النقرات على الغرسة مباشرة وإرسالها وحدها إلى الخارج: حيلة الضغط نفسها التي في العين.

\section*{ماذا تفعل نيورالينك}

تولّت شركة نيورالينك (\foreignlanguage{english}{Neuralink}) الجانب الهندسي من المسألة: خيوط رقيقة مرنة بدل الإبر الصلبة، وروبوت يغرسها، وغرسة لاسلكية محكمة الإغلاق. وفي عملها المنشور وُصف نظام موصول بالحاسوب بكابل، وذُكر الضغط والاتصال اللاسلكي بوصفهما خططًا. أما بيانات الإنسان المشلول اليدين الذي يتحكم في المؤشر فمعروفة حتى الآن أساسًا من تقارير الشركة نفسها؛ ووفقها في الغرسة نحو ألف قناة على بضع عشرات من الخيوط، وقد انزاح بعض الخيوط بعد الغرس. والفكرة نفسها ليست جديدة: مجموعات أكاديمية بمصفوفات أقطاب صلبة من نحو مئة قطب مكّنت الناس من التحكم في مؤشر، و”كتابة“ الحروف بقوة التفكير، وتحويل محاولات الكلام إلى نص على الشاشة بسرعة نحو ستين كلمة في الدقيقة.

\section*{القراءة والكتابة}

يقرأ المفكِّك ترددات النقرات لعصبونات كثيرة ويحصل على أقل من عشر بتات في الثانية: جزء ضئيل مما تحمله العصبونات نفسها. في المقابل يتعلم الدماغ والمفكِّك أحدهما من الآخر: يتكيف الإنسان مع المفكِّك، والمفكِّك مع الإنسان. متعلمان يتكيف كل منهما مع الآخر مسألة دقيقة: حتى حين يكون كل منهما مستقرًا وحده، قد يتأرجحان معًا.

الكتابة في الدماغ أصعب من القراءة. التيار عبر القطب يثير كل ما حوله، لا العصبون المطلوب وحده. يمكن ”رسم“ نقاط ضوء، كبكسلات مضيئة، فيميز الإنسان حروفًا بسيطة. لكن لا أحد يعرف حتى الآن كيف يكتب ترميز العين المضغوط بحيث يرى الدماغ صورة.

\section*{ما لا يراه المسبار}

القطب يسمع شبكة البيانات الهاتفية. أما المحطات الإذاعية، الدوبامين والسيروتونين والنورأدرينالين والأسيتيل كولين، فهي تراكيز مواد، والمسبار لا يسمعها. ولا يرى قوة الوصلات، حيث تُحفظ الذاكرة كلها. ولا يرى الألم كما يشعر به الإنسان. مصحح الأعطال الذي يرى البيانات ولا يرى السجلات سيبقى دائمًا مستغربًا لماذا يعطي المدخل نفسه أجوبة مختلفة.

\fullversion{17}
